% ================================================================
%  18.6 Les chantiers désolés (原书 p.894–908)
%  源文: _work/chunks/chunk131.md, chunk132.md, chunk133a.md, chunk133b.md
% ================================================================
\section{荒废的工地}

\subsection{(1) 悬而未决之事}   % 18.6.1

\notehead{176$'$}{}

\pmark\origpage{1194}(3月25日)昨夜,我在床上辗转了几个钟头,让自己重新浸入“Motif 的瑜伽”里,而不是像
本该的那样安然入睡。稍后,我也没有接着写笔记,而是又花了一两个钟头,潦潦草
草画出一堆蕴涵图,涉及我所知的、一个德拉姆上同调类(比方说,取特征零域上一个非奇异射影簇的上同调类)
要成为“代数的”所需的那些内蕴条件。我一共找到\textbf{十二}种路数,统统是 Hodge 猜想和
Tate 猜想的变体\fnote{971}{(*)}{(3月27日)这十二种变体,应当每一种都能对任意一个特征零的基本概形 $X$ 给出 $X$ 上一类相应的“系数范畴”(“系数类型”一词,取注“墓上的旋律——或自负”(n\textsuperscript{o}~167)中所说的含义)。倘若所论猜想为真,这条系数范畴应当把 $X$ 上的 Motif 范畴作为一个全子范畴(三角的)包含进来(这个猜想,在 $X$ 取一个域的谱这一特殊情形下,说的恰好就是这件事……)。细节我留给《数学反思》第3卷中专论 Motif 理论的部分(“Motif,吾爱”)。这也等于说,这十二种人人皆知的猜想的变体,给出(至少先天地)同样多种特征零域上“Motif”的不同观念。这就让将来会有十一位我朋友皮埃尔的摹仿者,各自“发现”属于自己的那种 Motif 观念,同时装作不知道别人的——尤其是(十五年来照例如此……)装作不知道某位已故者的,他主要正是以偏爱无用细节而著称……}。与此同时,我也说服了自己:要在一个 $\mathbb{Z}$ 上有限型的概形上定义“那个”(三角的)Motif 范畴——或者至少给出它的一个极紧的近似(假如这仍不是“那一个”对的)——所需的材料,应当差不多已经在手,只消再有一套“神秘函子”的理论即可;这个函子,是我在六十年代末曾公设过的\fnote{972}{(*)}{这个“神秘函子”的问题——它在特征 $p$ 的晶体上同调(经由“滤套 $F$-晶体”观念,$F$ 即“Frobenius”)与特征零的 $p$-进上同调之间架起“缺失的一环”——对我们的代数簇上同调认识而言显然至关重要,可将近二十年过去,自我以再明白不过的方式把它提出以来,这个问题至今无人认真问津……}。

在这里就此铺开,当然不合时宜。可我现在已经看清:正是在此时此刻,眼看 Motif 这个主题在我
把它留在一伙靠不住的人手里十五年之后,落到这般可怜的荒废地步,把当年我已抵达的那些思想
的几条主线勾画出来,恰是其时。我没有心思再等下去,等哪天有了闲暇(等《追寻栈》写完之
后)再去写那部本该写的“\textbf{那}部”系统著作;那份对一个梦的原原本本的讲述,本该作为让这
个\textbf{梦}终于扎根于精心酝酿成熟(并且发表出来……)的表述沃土之中的第一个\dmark{}\origpage{1195}大步骤,好让它
依着自己的本性开花结果。除了已为这部“数学虚构”之书预告并宣布过的第一块界石——即所谓
“六种运算”的代数对偶形式主义的一个概要——我还要在《数学反思》第3卷里附上一篇短文,提出
几个与 Motif 和代数闭链相关联的关键问题\fnote{973}{(*)}{照我目前的出版计划,《收获与播种》的前四部分(以\emph{葬礼}的第三部分亦即最后一部分作结)应构成《数学反思》的第1、2卷。第3卷将由《收获与播种》的第五部分(读 C.G.~Jung自传的笔记)和若干较短的文本组成,其中多数已在《引论》中预告。这样,《追寻栈》的第一卷就预定为《数学反思》的第四卷。}。眼看它们在坟墓里凋敝,我于心不忍;我巴不得看见它们重见天日,重新加入季
节的轮转……

思考回到\emph{葬礼}上来,已经五个多星期,而且再不肯离开。想必正因如此,“孤儿们”的念头——
他们被抛在一个病了的世界里,无人认领——近日又执拗地回到我心头。上一次有根有据地谈到其
中一位孤儿的笔记,是“墓上的旋律——或自负”(n\textsuperscript{o}~167),主题与昨夜以及稍前的这段
Motif 思考正相毗邻(我方才提到过它)。那正好是一个月前,就在我动笔写一条预定(这早已预先
定好)称作“四种运算”的笔记的前一天,我还不知道等着我的会是什么!到头来,一条笔记变成了
十六\fnote{974}{(*)}{(5月9日)写下这几行才不过两个星期,临头冒出的新事实,又把“四种运算”这场调查重新推动起来;它已经添了足足二十来条新注和子注。},我还以为永远做不完了——可到头来,我终究还是把这几件拖泥带水的“运算”走完了一遍!

而在这里,我最想做的,首先是回到那些孤儿身边,至少逐个叫出他们的名字——这也许对他们
有点好处,对我则一定有。第一次谈到他们,是一年前,就在那道同名的注“我的孤儿们”里,去年
3月底,与紧随其后的那道注“拒受遗产——或矛盾的代价”(注 n\textsuperscript{o}s~46、47)一口气写
成。写下这些笔记、给它们起这些名字时,仿佛有一股冥冥中的先知引路,我还没有料到,我留下
的那些东西,究竟“孤”到了何等地步\dmark{}\origpage{1196}——其含义之沉重、之切肤,连做梦我也想象不到;也
没有料到,我当时初次怯怯地指出的那道“矛盾”,会延伸到多远。这段记忆立刻又勾起另一段,是
前一个月的:我看见自己写下——仿佛是另一个比我更透彻的人借我的手在写——“\textbf{谁也斗不
过腐败}”。那是在写“无爱的世界”一节(n\textsuperscript{o}~19)的时候。我还记得,看见“\textbf{腐
败}”这个词白纸黑字地立在眼前,我起初愣住了。我心里某个“讲道理”的声音在训斥我:真是的,
你下手未免太重——“腐败”可是个大字眼,可不能信口开河!你还是换一副腔调吧!

我着实自我盘问了好一阵子,也许有几分钟。随后我明白,我既不会改掉这个“大”字,也不会加
一条注,解释说这个词是笔下一篇酣畅之际溜出来的,不必太当真。从那个世界里,这里那里,陆
陆续续回到我心头的这些“呛人的气息”,我心底有个人——比那个负责挑选“讲道理”标签的“我”
更有眼光——早就一清二楚它们意味着什么,根本不等我费神替自己编个说法\fnote{975}{(*)}{这段经过,我先是在去年3月“那张条子——或新伦理”一节(n\textsuperscript{o}~33)里讲过;两个月后,在发现\emph{葬礼}之后,又在 Notes 93--97 那一组原原本本得多的笔记里讲了,它们构成送葬行列之十“灵车”(与掘墓人同行)。}……

我也清楚记得那一天,那天的思考突然变了质地、我心中那\textbf{另一个}接手执笔的确切时刻。
那紧接在一段追忆之后:我数学生涯最初几年,数学圈子里那团裹着我的温情——多亏长者们给予
的接待,一直深到他们家里:Schwartz家、Dieudonné家、戈德芒家……转折就发生在我接着写下“显
然,对今天许多年轻数学家来说,要命的正是被切断了……与一切情感的潮流、与温暖的联系
……正是这一点折断了工作的翅膀,又夺去了它比一份郁郁寡欢、朝不保夕的营生更深的意义……”——就在
此刻,这个\textbf{无爱的世界}陡然在我眼前浮现,活了起来,再一次向我发问……

去年,“我的孤儿们”这个名字,是不请自来地落到我离去时留下的一切头上的(它们被我托付给
的那些近亲宣告了“死亡”……)。想必是因为,这个名字\dmark{}\origpage{1197}道出了一种\textbf{实实在在}、触手
可及的真实:我所“留下”或“托付”的,既不是“物件”,也不是“财产”,而是一些\textbf{活
的东西}。如今想起它们,我依旧把它们当作活的东西来想:茁壮、丰饶、生来就要生长、
绽放、又要孕育和催生别的茁壮而丰饶的活的东西。我确实感到自己留下了一份“财富”,可那不
是银行家的财富,而是园丁的、砌墙匠的财富:他们用双手让这些繁花似锦的园子、这些宽敞好
客的房子从地里长出来。这种贵重(乃至脆弱)之物的感觉,最紧地系于那些\textbf{观念}、那些
\textbf{问题}、那些我知道生机勃勃的\textbf{大主题}——我把它们留在了更年轻的人手里;这些东
西还需要劳作与呵护——远甚于那些我打磨到位的工具,或者那些我一路盖到头、布置停当的“房
子”\fnote{976}{(*)}{关于我心里那股催我“盖房子”(数学的房子)的冲动,见注“阴,作为仆人——与新主人”(n\textsuperscript{o}~135)。}。自有别人照自己的意思去那里忙碌,做他们自己的饭,享他们自己的清福;哪间房子嫌小
了,他们会照自己的需要扩建——就像我自己也常常不得不一扩再扩,虽然当年看上去,我“眼看
得很远”了。可是,我的数学家过去之所以仍旧攥得住我,靠的恰恰是\textbf{那些悬而未决之事},
是那些刚刚在绝好的地基上、用那么美的石料开了工的工地(而工人们已经走了,拿走了他们中
意的,糟蹋了剩下的……)。正是这些被撂荒的\textbf{工地}——今天我重访时,只见它们横遭劫掠、
残破不堪——我现在要把它们逐一巡视。

\subsection{(2) 吝啬鬼与老朽}   % 18.6.2. L'avare et le croulant

\notehead{177}{}

(3月27日)昨天一整天都交给了庶务。为了送去打字,我不得不把\emph{葬礼}第三部分亦即最后
部分的开头五十页重读一遍——连这里那里的小小措辞修整,外加又添上的几条脚注,足足花了五
个钟头。“阴阳之钥”的打字眼看就要完成。自从为这部分打字吃了那些没完没了的苦头之后\fnote{977}{(*)}{关于这些“苦头”(这还是委婉的说法),见注“祈祷与冲突”(n\textsuperscript{o}~161)的开头,以及《收获与播种》第五部分即最后一部分的开卷之注“Jung——或‘恶’与‘善’的循环”的开头。},我最终退而求其次\dmark{}\origpage{1198},找
了一位系里的秘书,在正职之外替我做。苦头到头了,谢天谢地——她做得又认真又麻利,一
个星期三十来页,无可挑剔。总能到终点的。也该到终点了!

除此之外,在一个 $\mathbb{Z}$ 上有限型的概形上,以三角范畴的形状构造 Motif 范畴这个问题,
继续在我脑子里打转——我又在床上熬了大半夜想它,没有睡觉——当心啊!起初似乎我的想法只
对特征零的概形行得通(比方说 $\mathbb{Q}$ 上有限型的),在基底 $\mathrm{Spec}(\mathbb{Z})$
本身上就不大灵;随后我记起,六十年代我就原则上确定过有限域上 Motif 范畴的结构。假定把我
当年做的工作明确化,我终于看到,一般情形下至少有一个完整描述的原理显现出来——说实话,
够紧凑的,但在我看来绝非不可企及。比起我六十年代的想法,唯一的新配料,是Mebkhout的哲学,
它凝结在那条说来奇特的“好上帝定理”里。除此之外,我拿来作假设性配料的,是“神秘函子”的理
论。眼下它还拿不出来,这肯定不是因为它“无从下手”(借一个我早领教过的说法\fnote{978}{(*)}{这正是我那位杰出的前学生Deligne用以埋葬“标准猜想”的斩钉截铁的考语——于是,在将近二十年里,我那些大胆的同时代人竟没有一个敢向它们下手。完整引文见注“悼词(1)——或恭维”(n\textsuperscript{o}~104)。}),而
是因为,我所认识的那些在代数簇上同调里做工的人,连在数学里也丢掉了对本质之物的分寸——
当然,他们为一场需要悉心照料的葬礼忙得不可开交……

平心而论,Deligne围绕Weil猜想做的工作——“Weil~I”,尤其是“Weil~II”——终归会在紧要关
头派上用场,到要为那些被指望表达 Motif 的系数范畴搭建六种运算的时候。尽管如此,还是得
有一位“糊涂”的、如人所愿地老朽不堪的已故者,在十五年之后灵光一现,从他那些亲爱的学生
与继承人乐于把他安置其中的软壁棺材里坐起来——而他本人对这一切毫不知情,而且差不多把生
前懂得的那一点点也忘光了——这样,在基底概形 $S$ 上描述 Motif 范畴的问题才算\textbf{一字一句地提
出}来,而搭造的原理(把与一个 Motif 相关联的一切已知结构要素都考虑在内的那种)也才顺带\dmark{}\origpage{1199}——纯属巧合地——终于得到明白的表述\fnote{979}{(*)}{正如我昨天在思考里预告的,我打算把这一描述收进《数学反思》的下一卷,附上一幅“Motif 广景图”的(极简略的)总轮廓——暗中的 Motif 把戏,闹得够久了。这里就预先声明:所设想的构造原理,不依赖代数闭链上任何一种猜想,诸如“Hodge”或“Tate”(或者昨天谈到的那十二种变体之一)。}。

1982年那部关于 Motif 的“难忘文集”之后,这份“Motif 藏宝”——十来年间一直是一个人的专
属禁地、一个秘密——似乎变成了三四个人共有的宝藏;这几位彼此往来,一派合谋者的神情,或
者像某个秘密教派的入会大佬。可是,把几个简单的问题白纸黑字写下来、提交给公众注意,不过
需要几天;要是想稍加用心地把它们界定清楚,明白交代手头有哪些配料、还须发展哪些配料,也
不过几个星期。从1970年起十五年,从“难忘文集”起三年,先是一个人,后来是那么几位,竟没
有一个肯从自己固然宝贵的时间里拿出这么几天——更不必说几个星期——这必定有极好的理由,
尽管他们当中没有谁打算去探个究竟。可他们煞费经营的那种气氛、他们赖以自持的那种心气,本
身就已经是一场发现之旅的堕落:发现之旅沦落为凌驾于他人之上、乃至蔑视他人的手段。这样的
气氛天然传播腐败,与创造背道而驰——哪怕乐在其中者是天才里最耀眼的几位。守着这样的心境
——一头孵着金币的守财奴的心境——他们切断了自身之中的创造之力,正如他们乐于扼杀他人身
上的这种力量。

\subsection{(3) 巡视工地——或工具与视野}   % 18.6.3. Le tour des chantiers - ou outils et vision

\notehead{178}{}

(3月30日)前天我满五十七岁,稍稍歇了歇工。只为“阴阳之钥”的结尾做了一点打字校对,昨天
又继续了一点。这是一桩让人歇息而愉快的活儿——至少,当打字的人自己也肯上心、一份我全身
心投入的稿子不至于被弄得面目全非回到我手里的时候,是如此。这回,我给自己放了两天假,\dmark{}\origpage{1200}细
细重读了五十来页清稿,就为了这里那里揪出一个还没放对地方的逗号……

不过,工作的劲头并不在顶点。几个星期以来,我心里有一种悲哀在提醒我:比起把手头这些笔记
一路写到它们自然的终点,还有更要紧的事在等我。我写着,像在逆水行舟;然而我知道,只要不出
意外、不遇不可抗之力,我只有给\emph{葬礼}终于画上那个句号,才会停笔。可把这股悲哀压下去、
放逐出去——它于是重得像一块石头——不让它在这些笔记里发声(除了眼下这一刻,顺带地、影
影绰绰地),这已经相当清楚地表明:我的思考,早在不知多久之前,就不再有资格称作“冥想”了。
它落在一种分裂之中:一边是写的人(却提防着不肯整个人扑进去\fnote{980}{(*)}{可是,就在上一段里,我刚写过(心里并无半点保留):托付去打字的那些稿子,我是“全身心投入”的。可见同样的话(或差不多的话……),随语境不同,可以有不同含义,或指示不同的细微差别。}),一边是活着、感受着的人(却也不停下来
“摆弄”他所经历的、浸透其中的意义)。在那里我感到,到达那个“句号”(当然,也不是说就把尚
待看见、尚待说出的这一点点草草了结……)然后返回我自身,已是刻不容缓……

除了笔记的活儿,这几天还有另一件事在给我调剂。那是一场数学思考,不以我意志为转移地重新
开了头。几天来我已明白:一套 Motif 理论形状的构造,带着二十年前我所望见的那全部气魄,离
“地平线”其实并不像我想的那么远。甚至有可能,一套“完全成熟”的理论——连六种运算的完整
形式主义(外加双对偶)在内——对肯把自己整个儿扑上去的人(不以掘墓的心性糟蹋自己的创造
之力)来说,只是几年功夫的问题。我还看出,要显式描述“那个”Motif 范畴\dmark{}\origpage{1201}——比方说,取绝对基底 $\mathbb{Z}$ 上有限型的概形(人们总应当能归约到这种情形)——有\textbf{两把}“钥匙”\fnote{981}{(*)}{其实还有第三把“钥匙”,这里不提,因为在我看来(不管对错)它没那么棘手。这就是“de~Rham--Mebkhout系数”的恰切定义(先不管滤套与 $F$-结构),比方说,取绝对基底 $\mathbb{Z}$ 上光滑概形。这一定义应当同时提供特征 $p>0$ 的一般晶体系数“那一个”恰切定义的钥匙——而我那些亲爱的前学生们(这一回首当其冲是Berthelot),始终没能、或者始终不愿把它提炼出来。
1983年6月(快两年了),Mebkhout向我讲解他围绕“好上帝定理”的“哲学”时,我的印象是:他对复数域 $\mathbb{C}$ 上光滑概形的(可构造的)离散系数范畴给出的“纯代数”(“德拉姆型”)描述,恰是Deligne在(别处已提过的)他1969/70年间在法国高等科学研究所所给讨论班(若无记错)中采用、却从未发表的那条进路的对偶;那条进路,靠的是带联络的 pro-模。我猜,从一个观点到另一个观点的过渡,经由对偶化函子 $\mathrm{R}\mathcal{H}om(-,\mathcal{O}_X)$ 完成——相对于所论概形的结构层——它把有限型的 $\mathcal{D}_X$-模(可以看作配备了一个可积联络的“ind-凝聚的 $\mathcal{O}_X$-模”)变成“pro-凝聚”的模(同样配备可积联络)。Mebkhout观点的长处,在于它为“好系数”给出一个简单而深刻的代数表达($M$-凝聚、完整性、正则性),这正是Deligne所缺的。Deligne观点的长处,在于它给出的是(与所要表达的超越性质的系数之间的)一个等价,而非反等价,并且更宜于表达所论系数范畴的乘法结构(张量积)。我猜,实践中人们会常常愿意在两张互为对偶的图式上同时工作。Deligne的解读,经由可积联络的模(或 pro-模)这一直觉,在我看来更贴近直接的几何直观。这一点的表现,尤其是(若基域为 $\mathbb{C}$):一个可构造的 $\mathbb{C}$-向量空间层,对应的是唯一一个带联络的 pro-模,而不是这种 pro-模的一个复形。正因如此(读者想必猜得到,这让我多么遗憾……),我预料最终最宜于用来发展六种运算的形式主义、并充当 Motif 范畴构造中第三种配料——钥匙的,恰恰会是他的观点——尽管正是他,毫无惋惜地把它埋葬了,仿佛借此把那位被拒不认的宗师所遗留的系数问题也一并埋掉……(5月9日)关于这一点,另见子注“……以及羁绊”(n\textsuperscript{o}~171(viii)),以及“五张照片”(n\textsuperscript{o}~171(ix))。}:一把,是“神秘函子”的理论,要有足够的普遍性与灵活性,好过渡到相应的三角范畴,把de~Rham--Mebkhout系数与寻常的 $p$-进系数(特征零)连起来;另一把,是在一个\textbf{有限}域 $k$ 上显式构造 Motif 范畴的问题(最好经由一种“纯代数”的构造,不援引 $k$ 上的代数几何),此外还有从 $k$ 上分离有限型概形(先从射影光滑概形开始)到这一范畴的“Motif 上同调”函子。后者我在当年已经构造到了差一个等价的程度,用的是Weil猜想与 Tate 猜想作启发\fnote{982}{(*)}{如果我记得不错,我当时只限于描述半单的 Motif 范畴。按同一原理,这一构造稍作变通,就能给出一个颇可信的候选者,充当不必半单的 Motif 范畴。我在这里说“Motif”,实际上说的是“iso-Motif”,即差一个同源不计的 Motif。但若对每个素数 $\ell$ 使用“$\ell$-进实现”函子,便能从那里重建不经同源的 Motif 范畴(其中诸 Hom 将是 $\mathbb{Z}$ 上、而非 $\mathbb{Q}$ 上有限型的模)。
说到我的构造“以启发的方式”使用 Tate 猜想,不可按字面理解。倘若确实存在(比如说,在有限域上)射影光滑概形上的上同调类,它们是“Motif 的”(这含义正待提炼)而非“代数的”(即并非来自代数闭链),那么 Tate 猜想(同理还有 Hodge 猜想,这回是在 $\mathbb{C}$ 上)就应当改述,把“代数类”换成“Motif 类”。假定人们果真能(如我下文所提示)在有限域 $k$ 上射影光滑概形的范畴上,定义出典范的(并且在适当意义上应当是“普遍的”)上同调函子,取值于已经构造好的范畴(所谓“$k$ 上半单 Motif 范畴”),那么这就顺理成章地给出一种形状完整的定义,界定人们将要称作“Motif 类”的上同调类:即 $\mathrm{Hom}(T^i,H^{\bullet}_{mot}(X))$(2i 维情形)的元素,其中 $T$ 是 Tate 对象,而 $H^{\bullet}_{mot}$ 是所拟的假设函子。这正是为什么,这个函子的构造,在我看来眼下是形塑 Motif 理论(不再是六十年代那种假设性的,而是形状完整的)的\textbf{头号}关键问题。}。

我对这一\dmark{}\origpage{1202}构造的正确性没有半点怀疑。尚待完成的工作——无疑明显更精细——在于:用给定的
有限域 $k$ 把这条范畴“钉死”,尤其是定义“Motif 上同调”函子,哪怕先只在 $k$ 上阿贝尔概形
的范畴上定义(这应当就足以把所求的范畴“钉死”……)。这第二个问题,在我看来,比起神秘函
子,较少技术性,而更直接是“几何的”。而且,它向我显形为\textbf{那把}能解开标准猜想\fnote{983}{(*)}{“标准猜想”一词在这里不可按字面取义,正如前一条脚注里的“Tate 猜想”。毋宁说,这些猜想的表述里,应当把所论的闭链类扩大(起初只限于代数闭链)。在“重新调整过”的标准猜想的“最终”表述里(即便它们按原样已然成立),“代数的”上同调类还要再换成“Motif 的”类。关于标准猜想,我将在“Motif,吾爱”(《数学反思》第3卷)里更详细地回头来谈。}的钥匙,并由此也解开特征 $p>0$ 的上同调理论里那些恼人已极的整性问题。凡此种种,使得这个问题对我有一股强大的吸引力!今晚已是第三个晚上,我重新拿起笔记,打算快速巡视一番那些在我看来最灼热的主题——十五年前我告别数学舞台时,被我的学生们、也被所有人抛下的那些主题\fnote{984}{(*)}{这些主题的一个极简略的初步“巡视”,见去年那道注“我的孤儿们”(n\textsuperscript{o}~45)。}。这一回,我总算要成功了!

\asterisms

\pmark\origpage{1203}\textbf{工地一:Topos}。我在这里提到它,主要是备而勿忘;在注“我的孤儿们”(n\textsuperscript{o}~46)
里,我已经相当周详地说过它。由于我的某些前学生——Deligne首当其冲——乐于糟蹋这条
关键性的统一性观念,自我离去之后,它被打入冷宫,只落得在边缘上苟活。如我在所引注里提醒
的,topos 与形形色色的“多重叠”其实在几何里俯拾皆是——不过,当然,尽可以看不见它们,正如
人们几千年里看不见对称群、看不见集合、看不见数零一样。

一套贴着拓扑直觉、柔软而精微的关于 topos 的语言,已在 SGA~4 的前两卷里被精心发展完备
(就是那块著名的“废话矿渣”——光辉的“SGA~4½”一书第一讲的引言里,Deligne如此称呼它)。这是
Leray所创“层”的观念周遭的语言与直觉的自然结晶;层直觉与层工具发展中的这第二步(或这“第
二口气”),在我看来,分量不下于第一步(其临时表达即戈德芒那本名著)。正是这一视野,让
$\ell$-进与晶体的上同调工具得以问世——然后,恰恰是那班佯装把工具据为己有的人,把它无限期
地埋葬了。

SGA~4 中关于 topos 的种种展开,并不自命完备与终极,但我认为,对 topos 式视野的大多数
直接的几何应用而言,它们绰绰有余。正如一般拓扑或寻常的层论,“topos 式的一般拓扑”本身,
在我看来并不提出什么真正深刻的问题。它是一套精心打造的语言,服务于对形状的拓扑-几何直
觉的某种扩展;这种扩展,是事物本身口授给我们的。这一视野长期蒙受的污名,以及落在它头
上的\textbf{嘲弄},在我眼里属于七十年代、八十年代数学世界的几大凋败之列。

这里说的,不是一间需要有人来救活的“破败工地”,而是一所全然竣工、设备齐全的房子:那些
住过它、又本被唤来在其中工作与生活的人,选择离它而去,还顺手撵走了盖房子的工匠。房子宽
敞、健全,一切各安其位,一如工匠出门做别的营生的那一天。它若还缺什么,缺的\dmark{}\origpage{1204}不是他这双手
的劳作,也不是任何人的劳作。也许,工匠本人对这些东西的敬意——那双手怀着爱做
出来的、他知道是美的东西——会让这股暴戾与轻蔑的气息消散,让这间本为迎接而造的房子重新
学会迎接。

\textbf{工地二:上同调的语言}。这首先是指导出范畴的语言;再者,是我为非交换上同调引入
的诸观点——两者都在五十年代下半叶。

头一条路线,本该由Verdier那篇著名的“论文”来承载;而Verdier本人对自己论文的埋葬
\fnote{985}{(*)}{关于这一点,见注“赊账的学位论文与全险”(n\textsuperscript{o}~81),以及“随心所欲的荣耀——或暧昧”(n\textsuperscript{o}~170(ii))。},同时也就是代数同调中导出范畴这一观点的埋葬。这一观点,曾在六十年代代数几何中上同调主题的繁盛期——尤其在对偶形式主义、不动点公式(Lefschetz--Verdier型)的发展中——扮演过关键角色。实践的需要,早已显出Verdier六十年代初发展的三角范畴框架之不足;这一框架,至今没有像本该的那样得到更新。

“非交换”一路,有Giraud的论文打下的良好基础工作,但它止步于一种 1-叠的形式主义,只能把
上同调对象直接几何化到二维为止。以 $n$-叠与 $n$-gerbe 的语言发展非交换上同调形式主义——众
多例子急切地提示了这一点——一路撞上严重的概念困难。鉴于在某个上流社会里,基础
问题跌入了人人不屑——说得更直白些,人人轻蔑——的境地,这些困难从来无人问津,直到两年
多以前,我亲自去啃\fnote{986}{(*)}{这指的是1983年2月我致Daniel~Quillen的信中的思考;在那封信里,我发现可以“一跃而过”那道豁开的“炼狱”——随着 $n$ 增大,($n$ 不严格的)$n$-范畴(或按我如今的说法,$n$-叠)的形状化描述中,似乎总要冒出来的、越收越紧的相容性关系。$n=2$ 的情形已然不是一桩轻松活,而 $n=3$ 的全部相容性,我想还没人有勇气把它们一一写出。这封信,如我下文还要提到的,成了长途跋涉《追寻栈》的发令枪;次月,乘着已经启动的思考的余势,我动身了。
这封信没有被视为值得一读,也没有得到回复。一年多之后,我才收到当事人的一点评语;对此我在“一段过去的重量”一节(n\textsuperscript{o}~50)里有交代。(参见打字稿第136页第二段。)}。

\pmark\origpage{1205}如今我看见两条路线在一门新学科里汇合——我在别处\fnote{987}{(*)}{见注“阴,这仆人——或慷慨”的子注 n\textsuperscript{o}~136$_{1}$(尤其打字稿第638页)。}提议把它叫作\textbf{拓扑代数}(l'algèbre~topologique):它是传统同调代数(当然是导出范畴风格)、同伦代数、(尚在娘胎里的)$n$-范畴、$n$-广群及相应的叠与 gerbe 的形式主义,以及 topos 的视野——它如今提供了已知最广的“纯代数”框架,供安放拓扑直觉——这几家的合成。这样一种合成的初始观念,六十年代已经齐备,其中包括\textbf{dérivateur}的观念——它注定要取三角范畴这个不够用的观念而代之,并且同样适用于“非加性”的语境。同伦代数中的某些重要进展,比如Bousfield与Kan七十年代初发展的同伦极限与同伦余极限诸观念——他们并不知道我的想法(在我亲爱的学生们那里,这些想法被当作“Grothendieck式的老糊涂涂鸦”打发了)——正落在这些想法的延长线上。

两年前,我以致Quillen的那封信\fnote{988}{(*)}{关于这封信,尤其见“一段过去的重量”一节(n\textsuperscript{o}~50,打字稿第136页第二段)。}起手,为我看见该做的工作起草了一份总施工图。那封信是《追寻栈》一书的发令枪;其第一卷(“模型的历史”)业已基本完成,无疑将作为《数学反思》第4卷出版。我预计还需要一两卷、一两年功夫,才能完成这一前期勘探——勘探的是一片极其丰饶的矿脉;二十年后,我似乎依然是唯一看见它的人。所以,这确实是一处撂荒了十五年的工地——但近一年来,它已在我手中复活。《概要》的写作,继之《收获与播种》的写作,曾把它打断;然而一旦《收获与播种》和那些(篇幅都有限的)文本——它们将与《收获与播种》最后一部分一起构成《数学反思》第3卷——写完,我就打算回去,把它做到底。

\pmark\origpage{1206}\textbf{工地三:六种运算,双对偶}。这指的是我引入Poincaré式或Serre式对偶形式主义(取离散或连
续系数)的那个观点。我当年起的“六种运算”这个名字,被我那班上同调学生细心地铲除
干净了。他们不过这里那里挑几样合用的来使,却把这几样合成一体所成的那种结构(连同双对偶
形式主义)扔进了盈亏账,尤其是,把这一观点所提供的那条无可替代的主线扔了(这条主线,对提
炼好的“系数范畴”至关紧要,见下文)。二十多年来,这套形式主义早已存在并经受了检验;可在“局
中”的人里,没有谁(除非是在注定保密、我无从知晓的文字里)费过神,去提炼那“代数总目”——
凡有这种“万能”对偶、以六种运算形式主义表达的地方,通用的那一套代数公式\fnote{989}{(*)}{(5月9日)在 SGA~5 最早的一次报告里,我曾煞费苦心,把这套公式说得详详细细;它将像发动机的神经一样,驱动整个即将展开的讨论班。这一讲是 SGA~5 里最关键的一讲,却在“屠杀版”里不见了踪影。全卷之中,竟无一字暗示它曾存在!见注“先祖”(n\textsuperscript{o}~171(i))里打字稿第942页脚注(*)。}。

可见,这里严格说不是一处“撂荒的工地”(要补的形式化工作微不足道),而是一个被系统性回避
的丰饶观点(正如 topos 的遭遇)。这场抛弃,无疑在很大程度上,造成了我在代数簇上同调主题上
看到的可悲停滞(少数例外\fnote{990}{(*)}{这些“少数例外”,主要是(1981年以前)Deligne那两项重要工作 Weil~I、II,以及晶体上同调、Barsotti--Tate 群的 Dieudonné 理论在特征 $p>0$ 的一般基底上的若干零星结果(后者是我在1969年前后开创的)。如我在别处指出的,在“好上帝-Mebkhout定理”的余波里,有过一阵回潮(定理与作者照旧一个被无视……),尤其是Mebkhout层理论的发展(它们被错误地叫作“反常层”,顶了本该属于别人的名分……),由Deligne等人展开。}除外)——尤其与我1955到1970年间赋予这一主题的强劲起飞相比。

如我已在《引论》(I~8,“一个秘密的终结”)里预告的,继《收获与播种》之后\fnote{991}{(*)}{提醒一下:这指的是《数学反思》第3卷,原则上它也收《收获与播种》的最后一部分。},我打算收入一篇短
文,勾勒“六种运算”形式主义的要点。多亏我那些\dmark{}\origpage{1207}学生的照料,今天它的存在本身已无人知晓——
除了曾直接参与 SGA~4(1963/64)与 SGA~5(1965/66)两个讨论班的人,而他们显然忘了\fnote{992}{(*)}{巧了,这也正是那篇以“核心”自居、名为(何其讽刺!)“SGA~4½”的文本劝人千万别去读的两个讨论班……(5月29日)六种运算视野的分量,见注“无用的细节……”(n\textsuperscript{o}~170(v)),部分(b)(“一些什么也干不了的机器……”)。}。这样,凡我力所能及的,我都做了:为一件效力完美的工具、一个在
上同调主题里不断把我们径直引向关键问题的丰饶观点,恢复名誉(只要还找得到守着好工具的工
匠)。

我刚巡视过的这三处“撂荒的工地”(或房子,或工具……),牵涉到的,更是一种\textbf{共同的代数
语言}——用来表达最形形色色的几何情形——而不是某个个别的几何情形,比如代数簇的上同调。
第二处工地,我称之为“拓扑代数”的,有时会让我跟一些无疑深刻的问题擦肩(比如球面同伦群
的某些问题),那是偶然,不是本意。我在那里的主要动机,一如既往,是发展出普遍性、灵活性
都足够的代数工具,来服务于那门尚在襁褓中的\textbf{算术几何}:从Weil猜想那颗胚胎起,我花了
一生中十五年的好时光孕育它、接生它、喂养它。几何实质的精华正在这门几何里;这些年里,它
才是我同数学的那场恋爱真正的核心——今天依然如此。接下来还要巡视的三处“最灼热”的主题,
说的正是这精华。

\textbf{工地四:“系数问题”}。这个问题,在Weil猜想的表述本身里,已然萌动\fnote{993}{(*)}{关于这一点,见注“操演”(n\textsuperscript{o}~169)的开头,那里我评说了Weil猜想的初始问题域。(5月29日)这一开头已独立成注“‘Weil猜想’的语境”(n\textsuperscript{o}~169~(i))。}。整个六十年代,它是我在上同调里的兴趣中心。对当时已约略望见的几种主要系数类
型,它已被提得清清楚楚,带着全部必要的普遍性与精确性\fnote{994}{(*)}{似乎并没有出现什么新式的“系数类型”,超出我六十年代下半叶所预见的那些。}。这一问题域——对理解代数簇上同调显然至关重
要\dmark{}\origpage{1208}——我早在第一次回顾自己的作品、那道以敬意命名的注“我的孤儿们”(n\textsuperscript{o}~46)里就说
过,又在注“墓上的旋律——或自负”(n\textsuperscript{o}~167)里回到它。两条根本的主线:一条,是刚才说
过的六种运算与双对偶的形式主义;另一条,是需要在一个或普或泛的基底概形上,给出基底域上
已知的诸种“系数”的合宜推广——它们(哪怕只是默而不宣)进入了对该域上射影光滑概形范畴已
知的那些上同调函子的描述:$\ell$-进上同调、晶体上同调、德拉姆上同调,乃至($k=\mathbb{C}$ 时)
Betti 上同调或 Hodge 上同调。

我不认为这样说有什么过分:这一问题域已经以萌芽的形式\fnote{995}{(*)}{作这一说明,我丝毫无意贬低Deligne与Mebkhout上述贡献的原创性与重要性,正如我不会因为说算术几何“早已在Weil猜想中萌芽”,就去贬低我自己的贡献在这门几何的诞生与最初起飞中的原创性与重要性。}包含着那套至今仍在等待问世的“完全成熟”的“Hodge--Deligne理论”,也包含着同样仍在等待的“德拉姆--Mebkhout系数理论”\fnote{996}{(*)}{大体可以说,先是Deligne的贡献(1969年前后),后是Mebkhout的贡献(1975年以后),分别朝着两个很不一样的方向,回应了定义合宜的“德拉姆系数”的问题——这种系数将能把光滑概形的寻常德拉姆上同调嵌进一个六种运算的形式主义。Deligne只在概形 $\mathrm{Spec}(\mathbb{C})$ 上定义了一类“好”系数,并且对 $\mathbb{C}$ 上分离有限型概形的结构态射 $X\to\mathrm{Spec}(\mathbb{C})$、对 $X$ 上的常值系数(可惜!)定义了函子 $Rf_!$ 与 $Rf_*$。Mebkhout定义了一类“好”系数,原则上对特征零域 $K$ 上任意分离有限型的 $X$ 都成立——但他没有进而对这类概形之间的态射 $f:X\to Y$ 定义 $Rf_!$ 与 $Rf_*$,也没有对 $Rf_!$ 与 $Lf^!$ 发展对偶定理(除了 $Y=\mathrm{Spec}(K)$ 的情形——而且即便那时,也只在复解析簇的超越语境里,那无疑明显更难)。Mebkhout迄今发展的理论(应当说,是在一种令人丧气到极点的氛围里发展的)还有一个局限:它目前只对光滑的 $X$ 做成(我猜,是因为没有系统地使用晶体观点;晶体观点为微分算子层提供了一个满意的替代物,而后者在光滑情形是那样便利)。};而这两套理论之所以至今仍在襁褓之中,没有像 $\ell$-进系数理论(素数 $\ell$ 与特征互素)那样,一两年间便臻于成年,原因是同一个:我那班上同调学生,Deligne首当其冲,忙不迭地把宗师遗留的问题域连同宗师本人一道埋葬。

\pmark\origpage{1209}Deligne这边(他享尽了科学界宠儿的一切便利),与Mebkhout那边(他身陷的孤立再彻底不过——
而这孤立,恰恰是最该接纳他的那些人压在他身上的),所迈出的步子固然零碎,却仍不失为宝贵的
引线,让人得以提炼出某些关键的系数范畴。写那道已引的注“墓畔的旋律”时,这些重要的贡献就
在我心里。此后,我又更深地浸入了那套六十年代已然成形的“系数与 Motif 的瑜伽”,如今有了一幅
更精确、也更完整的图景。因此,我打算在《数学反思》第3卷里,接着六种运算形式主义的概要,回头
再谈系数问题(同时也谈 Motif 问题)。

现在只需说:我分明看见,在一个或普或泛的基底概形 $X$ 上,有本质上的三种“基本”系数类型\fnote{997}{(*)}{我说“基本”的系数类型,是想借这称呼提示:我眼下能窥见的其余各种重要的系数类型,都应当能借这几种来描述——或以合宜的方式把它们“组合”起来,或对它们添加合宜的结构增益,或两者并举。在为de~Rham--Mebkhout系数所设想的诸种结构增益之中,有(除了似乎“内在于”所论系数范畴的“权滤过”之外)一种“德拉姆滤套”,它在 Motif 式的应用中扮演头等角色。有可能,这层附加结构(从六种运算形式主义的眼光看),唯有当它与一种贝蒂式的“离散”结构相结合时,才有意义;这应当使我们得以表述这一滤套所须满足的好性质。关于这些问题,我打算在“Motif,吾爱”(《数学反思》第3卷)里更周详地回头来谈。}:$\ell$-进系数($\ell$ 为\textbf{任意}素数);\textbf{de~Rham--Mebkhout系数}\fnote{998}{(*)}{提醒一句:对这种de~Rham--Mebkhout型系数,我如今看见互为对偶的两种变体:Mebkhout的那一种,与我迟疑着不愿称作“Deligne的”那一种——尽管后者是一个被他拒认的孩子!(5月29日)关于这个被拒认的孩子,见注“……以及羁绊”(n\textsuperscript{o}~171(viii))。关于德拉姆系数的细节,见注“五张照片(晶体与D-模)”(n\textsuperscript{o}~171(ix))。}(主要对在基底概形 $S$ 上有限型的 $X$ 有意义,最\dmark{}\origpage{1210}重要的情形是 $S$ 为环 $\mathbb{Z}$、$\mathbb{Q}$或 $\mathbb{C}$ 的谱);最后,是 Betti 系数(对 $\mathbb{C}$ 上有限型的 $X$)。这三种之中,唯有第三种,在我看来眼下已经确定下来,不带任何假设成分。要定义第一种(哪怕只对绝对基底 $\mathbb{Z}$ 上有限型的 $X$),或描述它与第二种
的关系,关键配料在我看来是那套“神秘函子”的理论——我六十年代末就公设过它,这个问题,看
来也同其余的一切一道沉了下去……到时候,我自当在合适的地方更周详地回到它。

\textbf{工地五:Motif}。关于我朋友Pierre~Deligne对 Motif 的埋葬——蒙教团全体祝福——我已
经说得相当周详,此处不必再度铺陈。这里我倒想强调一件刚刚才向我显形的新事,它本该在十五
年、二十年前就显形的。一个月前,在一个或普或泛的基底概形(比方说 $\mathbb{Z}$ 上有限型,或
只取一个代数闭域的谱……)上“形状完整地”构造 Motif 范畴,在我看来还是一件分明“在地平
线上”的事,淹没在遥远将来的迷雾里。这种心境,无疑是那些已然遥远的日子的顽固遗产:当年
Motif 的思考起步时,基地再假设不过,那时连 $\ell$-进上同调的形式主义都还没有。对我而
言还有一条“从轻发落的情节”:1958到1970年间,那些近在眼前的奠基性写作任务,把我的精力吸
干了;我的 Motif 思考(以及关于其他在当时那些紧迫任务面前显得“奢侈”的主题的思考),一直
被压缩到一份可怜的口粮,那份口粮,我几乎怀着逃学孩子的负咎才拨给自己!不管怎样,我留在
脑子里的印象是:系数问题,是马上就该做的成熟的事(不过得靠别人,因为我已经忙在别处
……);而 Motif,眼下顶多配上一本“数学虚构”的书,要是哪天我抽得出闲,去写一写。说实在的,
倘若我真的一头扎进去写它,而不是耗在那些后来世上没有一个人肯接手的任务上——尽管人人都
乐于享用我做的东西……——事情的步调,想必很快就会改观……

\pmark\origpage{1211}话虽如此,我终于看清了这件事——它本身,一旦摆在眼前,原是明摆着的:只要你肯花功夫去描
述足够“细”的系数,也就是说,照顾到一个 Motif 相关联的全部已知结构,你最终描述的,就是\textbf{Motif
本身}。或者说得对些:你最终描述的是一个范畴,它将把(三角的)Motif 范畴作为一
个\textbf{全子范畴}包含在内(这已经很可观了)——正如复数域上的 Motif 范畴(若承认 Hodge 猜
想一个相当强的版本)作为 Hodge--Deligne结构范畴的全子范畴而出现。至于要确切刻画——用与手
头系数直接相配的“代数”语言——这个全子范畴\textbf{究竟}是什么,也就是说,\textbf{哪些}系数恰
好“是 Motif”,那就落进了恐怕要棘手得多的问题。它们关乎一个 Motif 相关联的种种几何-算术
结构之间的\textbf{相容性}(这些相容性,我想我在前面引过的注“墓畔的旋律”里已经影射过)。解决
这些问题的路数(它们在我看来,与“Motif 理论”的实际构造毫不相干),兴许真的“要等一百年”。
不管怎样,经验一次又一次让我们看到:这类关于一个问题多么“无从下手”的预言,没有多大意
思,顶多是在勇气本没有挂稳的地方,劝人泄气……

(4月1日)关于“\textbf{Motif 的Galois群}(或\textbf{基本群})”的形式主义,再添几句评注。这个
概念(1964年,在我有幸结识我未来的前学生Pierre~Deligne之前,我已经提炼出来并开始发展)带
来一些极其精确、极其精微的直觉和形式主义。它的存在及其要点,不依赖于人们为一个域上
的 Motif(或任意概形上的“光滑”Motif)的观念所选定的特定构造——只要那个构造满足几条合
理的条件。我曾把这样一项任务交给Neantro~Saavedra:把我大约1964年提炼的辞典整理成可
发表的形状,放进尽可能一般的语境——辞典的一边,是我当年称作“刚张量范畴”的那些范畴里
的几何($k$-线性范畴,带满足若干合宜条件的“张量积”运算,$k$ 此处是一个\textbf{域});另一
边,是 $k$ 上 pro-代数群的线性表示理论(更确切、\dmark{}\origpage{1212}更一般地说,是 $k$ 上的“pro-代数 gerbe”的表
示理论)。他把这项任务在他的论文里做到底,1972年由“数学讲义”出版(LN~265)
\fnote{999}{(*)}{(5月10日)这几行写下之后,我有机会读到了那本书——作者可没觉得有必要寄我一本。我看清了:书中,Saavedra俨然是那里所展现的新哲学的光辉发明人;他一丝不苟地遵照我交给他的那些笔记,而几乎不提我的名字(无论是书中引入的观念与关键结果,还是一些早已为人所知的观念,比如晶体、分层模或 Motif)。他为那个主要观念重新起的“塔纳卡范畴”这个名字,是一场天才绝伦的骗局——那名字,断不会比他所自命的这套理论,更是他自己的发明。何况这桩“父亲身份”本来就是暂时的:我的朋友皮埃尔已经在丛书问世十年之后,把必要的事情料理妥当,好叫它(如众人所期待的那样)回到那位早已内定的父亲名下。这场对一具尸身的辉煌手术的细节(在“SGA~4½--SGA~5”那场以同样不可模仿的手法进行的手术之前,它是第一场、也是唯一一场这种规模的),见那组注“第六根钉子(钉进棺材)”(n\textsuperscript{o}s~176$_{1}$至176$_{7}$)。}。这本辞典我当年还推进得更远(尤其涉及把滤套结构或分次结构等,翻译到某些纤维函子上,或涉及与塔纳卡范畴相关联的一个“极化”观念),这些在Saavedra的论文里都没有做\fnote{1000}{(*)}{(5月10日)这原来是一句错误的揣测。它出自我的一个信念:Saavedra断无能力“收尾”我给他指的那个计划——单是“pro-代数 gerbe 的线性表示”这一观点的驾驭,很长时间里似乎就超出了他;他的数学行囊又极其单薄。以Saavedra那点并不出众的本钱,从我离去(那时他对上同调、对代数群的构造一无所知)到书出版,不足两年,他竟有机会消化(而且消化得尽善尽美——书的成色即是明证)书中随手把玩的那一大堆四面八方的观念——这在我看来是不可设想的。关于这一点,见那组已引注“第六根钉子(钉进棺材)”里的注“Monsieur~Verdoux——或献殷勤的骑士”。},或者在“难忘文集”LN~900里都没有(在那里,Saavedra
的论文被重做了一遍,而 Motif Galois群的观念居于问题的中心——只是我的名字,在这个话题
上,同在所有其他关于 Motif 的话题上一样,一次也没有被提起)。

我还想指出:前面谈到过的那第一步——确定有限域上 Motif 范畴(差一个等价)——原本是对
该有限域的 Motif\dmark{}\origpage{1213} Galois群的确定\fnote{1001}{(*)}{(5月10日)我注意到,这一确定,也见于Saavedra那本取之不尽的书里(对我这卑不足道的人只字不提,这还用说么)。它用的是整体类域论的上同调理论(确定群 $H^2(\mathbb{Q},T)$,其中 $T$ 是 $\mathbb{Q}$ 上一个乘法型群)——所以,这也算我的前学生(其本钱俨然超乎常人)不到两年里就消化掉的东西之一……},该群应当是交换的(它被 Frobenius 元素拓扑地生成),实际
上是一个经某些 $\mathbb{Q}$ 上的 pro-代数环面扩张的 $\widehat{\mathbb{Z}}$(由 Frobenius 生成)\fnote{1002}{(*)}{这里说的是对\textbf{半单}Motif 分类的那一个 Motif Galois群。要得到一般的 Motif,还须用 $\mathbb{Q}$ 上的加法群 $\mathbb{G}_a$ 作它的乘积。}。第二步,是 $H^2(\mathbb{Q},T)$ 中那个元素的描述——按照Giraud的理论,它恰把纤维函
子的 $G$-gerbe 加以分类\fnote{1003}{(*)}{要点在于:这个类在所有 $\ell\neq p=\mathrm{car.}\,k$ 的位置上都变为零(多亏“$\ell$-进上同调”纤维函子的存在),而晶体纤维函子的存在,让我们对这一类在缺失的那个位置 $p$ 上的命运,有了足够的信息。}。

正如我在注“一个梦的回忆——或Motif的诞生……”(n\textsuperscript{o}~51)里所说,我是在寻找下述联系时撞
见 Motif Galois群的:对一个投射有限Galois群 $\mathrm{Gal}(\overline{K}/K)$,取变动的素数 $\ell$ 的各
个 $\ell$-进表示——例如取 $H^i(\overline{X},\mathbb{Q}_\ell)$ 得到的那些,其中 $X$ 是射影光滑概形,$i$ 为整
数(或它的某个合宜的子模)。Serre看的是Galois群在 $\mathrm{Aut}(V(\ell))$ 中的像——那是一个约化的
$\ell$-进李群——而这个像的结构(按李理论的意义)似乎与 $\ell$ 无关。正是在探寻这一现象(直到今
天,它本身仍属假设)的深层原因、把它与 Tate 猜想相联系的过程中,我发现了 Motif Galois群的
观念——紧接着“Motif”与“Motif 上同调”的观念之后。

如果说我曾把一件既简单又深刻的东西带到世上,如果说我的数学家生涯里有过创造之举,那正
是这条沟通几何与算术的关键观念的诞生。也正因为如此,去年那个难忘的4月19日,当看到这件东
西被人以那样一副气定神闲的做派据为己有——仿佛它是某人在一个技术段落的拐弯处,即兴哼出
的最末一句小玩意:瞧,简单得像剥白菜,只消把我们的那篇小文章里陈述塔纳卡范畴理论的命题
4.7.3 拿来用在这里就行——我给一股难以想象的\textbf{无耻}之气噎住了\fnote{1004}{(*)}{写这几行时,又逼着我想起一件十分相似的事:把一个 $\ell$-进层系数的 $L$-函数的定义引入,同样不提任何人,同样仿佛是信手拈来的最末一条老生常谈,由同一位光辉的作者即兴完成。关于这一点,见注“操演”的子注“……以及‘废话’”(n\textsuperscript{o}~169$_6$),打字稿第891页。}。七十年代里那些光辉的\dmark{}\origpage{1214}先例之后\fnote{1005}{(*)}{甚至早在六十年代——见“驱逐”一注(n\textsuperscript{o}~63)。},八十年代的数学,就是这样做的。

不过我扯远了,好吧——我本该是领人参观一处工地的,不是来抒情的。那就说明一下:同投射有限
基本群的情形一样,若 $X$ 是域 $k$ 上几何连通的概形,应当区分概形 $X$ 本身的 Motif 基本群,与
“\textbf{几何的}”Motif 基本群。两者并不重合,\textbf{哪怕} $k$ 代数闭——因为 $k$ 的 Motif 基本群
并不因此就是平凡的(它连通,仅此而已!)。所以要引入 $X$ 的“几何的”Motif 基本群;它的职责
(除了别的之外)是在与几何投射有限基本群 $\pi_1(\overline{X})$ 相关联的诸 $\ell$-进李群(作为其商)之间
建立联系。它的定义,是自然同态
$$\pi_1^{mot}(X)\to\pi_1^{mot}(\mathrm{Spec}(k))$$
的核(相对于 $X$ 上光滑 Motif 范畴上一个纤维函子的选取)。

我想要抵达的地方正在这里:这个核——不妨记作 $\pi_1^{mot}(X/k)$——应当是构造“$X$ 在 $k$ 上的
(几何的)Motif 同伦型”的第一步;这一同伦型,我前文已经顺带影射过\fnote{1006}{(*)}{见注“为一副模糊的骸骨作安魂曲”(n\textsuperscript{o}~165)。}。“这个同伦型”的形状完整的描述\fnote{1007}{(*)}{作为对象的类型,我预料它将是一个相对同伦型(按Illusie的意义),放在 $\mathrm{Spec}(K)$ 的 fpqc topos 的“扩张”topos(按Giraud的意义)里——该扩张对应着 $X$ 上光滑 Motif 范畴的纤维函子的(在这个 fpqc topos 上的)gerbe。这个同伦型在(刚才描述的)基底 topos 上的相对上同调,是拟凝聚的(甚至是“凝聚的”),并且可以等同于 $X$ 在 $K$ 上的 Motif 上同调。若用 $X$ 的一个复点(特征零的 $K$),取 Betti 纤维函子,则相应的纤维同伦型应当典范地同构于经超越途径与 $X\otimes_K\mathbb{C}$ 相关联的 $\mathbb{Q}$-同伦型(不计挠现象……)——至少当 $X\otimes_K\mathbb{C}$ 是 1-连通的时候。},它
的“上同调”应当不是别的,正是 $X$ 的 Motif 上同调;这属于“Motif”工地上一个有意思的概念性
工作,朝着一条分明不同的方向\dmark{}\origpage{1215}(而且很大程度上多半独立)——不同于那项中心任务:
Motif 范畴的实际构造,以及为它们建立六种运算的形式主义。

\textbf{工地六:标准猜想。}如我在前面一条脚注(打字稿第1202页脚注(*))里解释过的,这些猜
想可以在两种不同的意义上理解。首先是字面的意义,如我1967年在孟买讨论班上陈述它们的那样
\fnote{1008}{(*)}{Algebraic~Geometry,~Bombay~1968,~Oxford~University~Press~(1969)。}。照这个样子,在我看来,它们概括了代数闭链理论眼下最关键的问题——至少从这些闭链的所谓“同
调”等价的观点看是这样。

陈述这些猜想的时候,我的主要动机,其实并非朝向闭链本身,而是朝向它们(或许……)提供的
手段:在一域上建立一套半单 Motif 的理论,满足那些本该在十五年、二十年前就“人所共知”(至今
却仍讳莫如深……)的要求。我将在《数学反思》第3卷里指出这些猜想的种种减弱的变体——它们足
以建成这样一套理论(其中最弱的那个,几乎就是为了这个目的而必要且充分的)。如我在别处强
调过的:即便初始形态的猜想被证明在某个给定的域 $k$ 上成立(比如 $k$ 有限,甚至一切 $k$),这也
绝不意味着,应当称作“Motif 的”上同调类\fnote{1009}{(*)}{我想我可以提出一个合理的定义,界定射影光滑代数簇上的 Motif 上同调类——至少当基域特征零时如此。对一般情形,关键的(前文已经谈过的)是有限基域的情形。知道了后者之中 Motif 类的描述(且听下文分解),我想就能给出“那一个”Motif 类的恰切定义。可参照打字稿第1202页脚注(*)处的评注。}(人们可以指望这些类使 Hodge 猜想、Tate 猜想之类种种猜想
为真)就必然是代数的。倘若有一天发现存在非代数的 Motif 上同调类,那多半意味着:代数闭链在
Motif 理论中——也就是在代数簇上同调的算术-几何研究中——的分量,要比理论初创时我有理
由相信的更轻一些。话虽如此,我如今隐约望见的 Motif 理论的实际构造,先天地独立于关于
代数闭链的时下诸猜想(Hodge、Tate 或“标准”型)。

\pmark\origpage{1216}但这并不妨碍:一方面,标准猜想及其种种变体;另一方面,Hodge、Tate 猜想及其众多的变体
——这些猜想,尤其包含代数闭链的\textbf{存在性}命题(即上同调类的代数性),或者(在修改版
里)所谓“Motif 的”上同调类的存在性命题——彼此紧密相联,并且同主要“系数类型”的描述、乃
至同 Motif 范畴本身的描述,紧紧相连\fnote{1010}{(*)}{这并不与我刚才的断言相抵触——我说过,我隐约望见的 Motif 范畴(比方说一个域上的)构造,与所论的种种猜想“独立”(即“技术上”或“逻辑上”的独立)。我说的那些“紧密的联系”(比如,我在 Hodge、Tate 型猜想中看到的十二种变体,恰好提示了同样多的不同类型的上同调“系数”),属于启发而非技术——正如(所谓“猜想性的”)Lefschetz--Verdier公式与 Frobenius 对应的迹公式之间的联系一样。在后一情形,这个根本的启发式联系(并非逻辑依赖的联系),已在注“操演”的两条子注“\textbf{真正的}数学……”、“……以及‘废话’”(n\textsuperscript{o}s~169$_5$、169$_6$)里得到郑重的指明。}。

这里,又是一件近二十年来早该做的澄清、归置、通报的活儿,没有被做(尤其,没有公之于众),
做的人却宁愿到今天还在埋葬丰饶的观念(当它们还不曾发表),或者拆毁它们(当它们发表了),
好独占(眼前的)利益与(日后的)功劳,而不肯把这些迷人的问题通报出来、交到大家手里——
这些问题,对我们理解几何、拓扑与算术之间的纽带,至关紧要。我看出,这里欠缺的,绝不是能
力,甚至不是光辉的天分,而是一种朴素的诚实,以及在对一个派发声望与权力的“科学共同体”的
关系里,某种起码的\textbf{廉耻}——在那些并不因此觉得自己欠下丝毫义务、欠下丝毫“服务”式“回
报”的人身上。所以,虽然我同这个题目断了往来已十五年以上,虽然我已经说不上“局中”二字,
费力重新入水的却还是我:至少要在《数学反思》第3卷里,把我那些更年轻、更有天分的人漏下的,
尽力修补;到底,把他们没有慷慨去做的事做了。

到这里,我相信我已经把这些“工地”巡视了一遍——从建造那门“算术几何”的眼光看(这门几何
的基石,是我整个六十年代铺下的),它们如今(而且自我告别数学舞台那一刻起)正是“\dmark{}\origpage{1217}最灼热”
的。我绝无意说,我已经把那些实质问题——那些也许只有我看得见、而又放在我心
上的问题——粗粗走遍了。据我所知,它们仍停留在我离去时的原地,其中许多甚至还没来得及在
文献里得到显式的陈述。其中我要特别指出\textbf{离散Riemann--Roch猜想}(概形框架下的)\fnote{1011}{(*)}{这一猜想似乎第一次被显式陈述,是在那道名字意味深长的注“屠杀”的子注 n\textsuperscript{o}~87$_1$ 里——因为这条猜想,本就是 SGA~5 遭屠杀的诸事项之一,在Illusie的编订版里消失得连一个\textbf{名字}都没留下。}。同样,还有把\textbf{局部与几何整体类域论}推广为一条本质
上“几何”性质的\textbf{对偶}命题(经典“算术”命题作为推论给出)的工作。1976年致Larry~Breen的几封信里谈到了它,那几封信将作为附录收入《追寻栈》第一章(因此会出在
《数学反思》第4卷)。在这些命题里,主要的工作将在于对人们所处理的“系数”范畴的细心描述。其中
一个重要角色,由Serre发现的一种自对偶性扮演\fnote{1012}{(*)}{除了Serre这个漂亮的念头之外,影响我的,还有 Lang 在几何整体类域论中、Serre在局部类域论中引入的“几何”观点。}:特征 $p>0$ 的域 $k$ 上、幂幺代数群范畴里(差根同源不计)的一种自对偶(在我看来,除了当年听我讲过的屈指可数的几个人,至今无人知晓)。把这类命题推广到更高维的问题,(至少对我)是个彻底的谜(不过Milne在代数曲面的情形,或许有些眉目……)。

这些对偶问题,我记得可以上溯到五十年代末。那时我还顺手开创了一个与之相联的构造:与一个
域上有限型概形(先从这里开始……)相关联的、pro-代数群的“\textbf{雅可比}”(链)\textbf{复形}——
借与那些局部环相关联的合宜的“局部雅可比簇”,类比于我几年前在凝聚对偶里构造过的“残差”
或“对偶化”复形。所有这些对偶问题,在六十年代,都被推到了第二线,尤其是为了
\dmark{}\origpage{1218}发展平展与 $\ell$-进上同调的“废话”和 topos 的语言。我纲领里相对局部与整体雅可比的那一小段,约在1977年由Contou-Carrère完成(没有提我这卑不足道的人);而一见到Deligne与Raynaud对它摆出的脸色,他就赶忙收了起来\fnote{1013}{(*)}{见注“棺材3——或有点太相对的Jacobi簇”(n\textsuperscript{o}~95),及其子注 n\textsuperscript{o}~95$^{1}$,其中谈到Contou-Carrère在数学大世界里的几段倒霉遭遇。}。如今,要重新拾起、发展那些打着我太明显印记的观念(哪怕当事人力图把这印记
藏得最好),是需要一点勇气的。唯一死磕到底的,是Zoghman~Mebkhout;而他遭到的命运——在反常
层会议上登峰造极的那些“绝技”——足够清楚地显示了这里头的风险。

倘若我想要开列一张1955到1970年间我发现的漂亮问题的清单(它们当中的一些,当年我曾在周
围这里那里谈论过),我还得再花上几天;若想稍稍说得明白些,把来龙去脉都交代了,恐怕要几
个星期。这里不是做这件事的地方,我也怀疑我永远会去做。更何况,倘若我盼着有一天(谁知道
呢!)哪位年轻数学家肯在其中一个问题上扑进去,练练手,再让自己为人所知,那还是让他自己
重新发现的好——犯不上冒险,让人给贴上某种标签。

当心将来等着我们的那些反常层会议……
